\begin{textsample}{POS Dim 3 – ap – Score 14.00 – ap/ap\_inf\_17.txt}  \label{ex:f3_pos_145}
3.6.3.
Etapas de Transição A \textbf{criança} da Educação \textbf{Infantil} passa por momentos distintos de transições que são : transição da casa para a escola, transição no interior da instituição e a transição da Educação \textbf{Infantil} para o Ensino Fundamental ( BRASIL, 2009 ).
É necessário que esses momentos sejam pensados, planejados e muito bem organizados para que a \textbf{criança} se sinta segura e \textbf{acolhidas} nesse processo.
Ao fazer a transição da casa para a instituição de Educação \textbf{Infantil}, a \textbf{criança} sai de sua zona de \textbf{conforto}, pois deixa de conviver somente no vinculo familiar, com pessoas, costumes e regras conhecidas, para um espaço totalmente novo, ambiente coletivo, com regras e pessoas inicialmente estranhas.
Nesta fase, a segurança emocional dos pais e as formas de adaptações/acolhimento por parte da instituição são fatores determinantes para o bem-estar da \textbf{criança}.
É importante que os pais estejam seguros em deixar seus filhos na instituição, fazendo com que a \textbf{criança} se sinta segura, e ainda, que a escola seja um espaço que transmita aconchego e encantamento.
A adaptação ou acolhimento da \textbf{criança} é um processo lento e gradual, em um tempo indeterminado, com avanços e retrocessos, sendo essencial o respeito ao tempo de cada \textbf{criança} e de sua família, haja vista, que esse processo é encarado de maneiras diferentes pela \textbf{criança} ( BRASIL, 2009 ).
Algumas \textbf{crianças} podem ver a escola como um lugar seguro, divertido e que rapidamente criam vínculos com a \textbf{professora} e com as outras \textbf{crianças}, no entanto, irá acontecer também de ter \textbf{crianças} com sentimentos oposto, que vão sentir -se com medo, angustiados e irão sofrer com a separação da família.
Com base nas DCNEI/2009 recomenda -se as seguintes ações para a efetivação da adaptação e acolhimento : - Planejamento : é importante que a escola considere todos os aspectos desse período de adaptação, pensar em tempos, espaço acolhedores, materiais didáticos que desperte a \textbf{curiosidade} e interesse da \textbf{criança}, e distribuir atribuições para todos os profissionais da escola. - Envolvimento dos funcionários da escola : todos da escola independente de sua função é (co)responsável por esse processo de acolhimento da \textbf{criança}.
Para isso é importante que a escola promova palestras e cursos com toda equipe escolar sobre a importância do acolhimento.
Assim, estarão preparados para esse momento, entendendo suas atribuições e importância para a qualificação da chegada e permanência da \textbf{criança} no contexto escolar. - Participação das famílias : é essencial para esse processo o envolvimento da família, efetivando uma parceria de \textbf{cuidado} e educação, criando um vínculo de confiança.
Para isso, é necessário, que a família conheça o papel da escola, seu funcionamento e seus profissionais. - Atendimento à diversidade : é necessário que a escola esteja preparada para as manifestações individuais das \textbf{crianças}, suas particularidades e especificidades. - Lidar com os sentimentos : no período de adaptação acontece uma explosão de sentimentos, seja por parte das famílias que ficam inseguros em deixar seus filhos, ou pelo profissional que precisam lidar com reações diversas por parte das \textbf{crianças}, como : choro, inquietação, recusa de alimentos entre outras.
É nesse momento que a escola intervém com paciência para aproximar as \textbf{crianças} da \textbf{rotina} da escola, criando vinculo e confiança com as \textbf{crianças} e entre seus pares.
É importante destacar que essa adaptação também precisa ser considerada quando ocorre mudança de instituição escolar ou entre as etapas da Educação \textbf{Infantil}.
Essa mudança de uma turma para outra dentro das etapas da Educação \textbf{Infantil} é outro processo de transição que também precisar ser bem trabalhada.
Ainda que todas as turmas façam parte da mesma etapa, possuem mudança em suas \textbf{rotinas}, de pessoas, de espaços.
A preocupação de criar condições favoráveis para a transição entre turmas requer atenção de todos os \textbf{adultos} que cuidam da educação da \textbf{criança}.
A nova estrutura que compõe a Educação \textbf{Infantil}, com base na BNCC ( BRASIL, 2017 ), é um fator que favorece esse processo de transição, pois as aprendizagens essenciais estão organizadas por campos de experiências, seus objetivos de desenvolvimento são estabelecidos gradualmente, ou seja, é possível perceber a evolução em cada faixa etária, com isso é possível à continuidade do processo educativo.
Em relação a transição entre as etapas a Educação \textbf{Infantil} e Ensino Fundamental.
As DCNEI recomendam : > Art. 11.
Na transição para o Ensino Fundamental, a proposta pedagógica deve prever formas para garantir a continuidade no processo de aprendizagem e desenvolvimento das \textbf{crianças}, respeitando as especificidades etárias, sem antecipação de conteúdos que serão trabalhados no Ensino Fundamental ( BRASIL, 2009, p. 5 ).
Para a continuidade do processo educativo, torna -se necessário que as escolas considerem as orientações do Parecer n. 20 - CNE/CEB : > [... ] c ) planejar o trabalho pedagógico reunindo as equipes da \textbf{creche} e da \textbf{pré-escola}, acompanhado de relatórios descritivos das turmas e das \textbf{crianças}, suas vivências, conquistas e planos, de modo a dar continuidade a seu processo de aprendizagem ; d ) prever formas de articulação entre os docentes da Educação \textbf{Infantil} e do Ensino Fundamental ( encontros, visitas, reuniões ) e providenciar instrumentos de \textbf{registro} – portfólios de turmas, relatórios de avaliação do trabalho pedagógico, documentação da frequência e das realizações alcançadas pelas \textbf{crianças} – que permitam aos docentes do Ensino Fundamental conhecer os processos de aprendizagem vivenciados na Educação \textbf{Infantil}, em especial na \textbf{pré-escola} e as condições em que eles se deram, independentemente dessa transição ser feita no interior de uma mesma instituição ou entre instituições, para assegurar às \textbf{crianças} a continuidade de seus processos peculiares de desenvolvimento e a concretização de seu direito à educação ( BRASIL, 2009, p. 17 ).
Ao considerar os dois documentos é evidente a necessidade do trabalho conjunto entre as etapas, para que haja um \textbf{equilíbrio} no momento da transição, de forma que a etapa seguinte considere os que os educandos sabem e são capazes de fazer, garantindo a continuidade do trabalho pedagógico.
Por esse motivo a importância do acesso aos relatórios, portfólios ou qualquer outro tipo de \textbf{registro} de \textbf{acompanhamento} da \textbf{criança} de sua trajetória na Educação \textbf{Infantil}.

% matched lemmas: acolher, acompanhamento, adulto, conforto, creche, criança, cuidado, curiosidade, equilíbrio, infantil, professora, pré-escola, registro, rotina
\end{textsample}
